\documentclass[10pt,dvipdfmx]{jsarticle}
\usepackage[width=160mm,height=230mm]{geometry}
\usepackage{ascmac}
\usepackage{graphicx}
\usepackage{enumerate}
\usepackage{url}
\usepackage[newfloat]{minted}   % requires Pygments---a Python package
\usepackage{listings,jlisting}
\renewcommand{\lstlistingname}{ソース}
\lstset{language=sh,%
        basicstyle=\footnotesize,%
	commentstyle=\textit,%
	keywordstyle=\bfseries,%
	frame=shadowbox,framesep=4pt,%
	showstringspaces=false,%
	numbers=left,stepnumber=1,numberstyle=\footnotesize,%
	numbersep=1zw,%
	lineskip=-0.5ex,%
	xleftmargin=2zw%
	}%
\usepackage[lighttt]{lmodern}
\input{listings-javascript.prf}

\pagestyle{myheadings}
\newcommand{\maru}[1]{{\rm\ooalign{\hfil\lower.168ex\hbox{#1}\hfil\crcr\mathhexbox20D}}}%
\newcommand{\answer}[1]{\begin{flushleft}
            \framebox{\large\bf #1}\end{flushleft}}


%
%%%%%%%%%%%%%%%%%%%%%%%%%%%%%%%%%%%%%%%%%%%%%%%%%%%%%%%%%%%
%
%  解答に際しての注意
%
%%%%%%%%%%%%%%%%%%%%%%%%%%%%%%%%%%%%%%%%%%%%%%%%%%%%%%%%%%%
% 1 試験時間は09:00-12:00である．
%
% 2 採点にあたっては，内容はもちろんであるが,
%   出力された文書の体裁も評価する．
%
% 3 問題文を削除しないこと．
%
% 4 解答はファイル名「exam-2022.pdf」としたものをTeamsで提出すること．
%
% 5 解答にあたっては，参考資料等を参照してもよい．
%
% 6 他人との相談も差し支えない．
%   ただし，コミュニケーションは，コンピュータを介したものの
%   みを認める．つまり，相手の感覚器に直接訴える音声や画像，
%   触覚による直接的な情報交換，携帯電話の使用などは固く禁じる．
%   教員・TAへの質問はこの限りではない．
%
% 7 明らかなファイルのコピーが認められた場合，コピーした者，
%   コピーされた者にはペナルティを課す．従って，解答ファイル
%   のパーミッションを適切に設定(chmod 600推奨)することも重要である．
%
% 8 解答を記述する前にこのファイルをコンパイル & プレビューし，様子を確認
%   しておくこと．
%
%%%%%%%%%%%%%%%%%%%%%%%%%%%%%%%%%%%%%%%%%%%%%%%%%%%%%%%%%%%
%
%
%　☆重要☆　必ず変更してください
%
% 「愛地　博」の部分　⇒　自分の名前に
% 「20yy3110xx」の部分　⇒　自分の学籍番号に 
%  
%
%%%%%%%%%%%%%%%%%%%%%%%%%%%%%%%%%%%%%%%%%%%%%%%%%%%%%%%%%%%
%
\newcommand{\氏名}{神野　友哉}  % {}内を自分の氏名に変更してください
\newcommand{\学籍番号}{2022311042}  % {}内を自分の学籍番号に変更してください
%
%
%
%
%%%%%%%%%%%%%%%%%%%%%%%%%%%%%%%%%%%%%%%%%%%%%%%%%%%%%%%%%%%
% 以下の設定には手を加えないこと - ここから
%%%%%%%%%%%%%%%%%%%%%%%%%%%%%%%%%%%%%%%%%%%%%%%%%%%%%%%%%%%
\markboth{{\氏名}({\学籍番号})}{{\氏名}({\学籍番号})}
\title{\framebox{\hspace{4zw}2022年度コンピュータリテラシ定期試験%
                 \hspace{4zw}}}
\author{\氏名 \\
        \学籍番号 \\
        愛知県立大学 情報科学部 情報科学科}
\date{2022年8月3日}
%%%%%%%%%%%%%%%%%%%%%%%%%%%%%%%%%%%%%%%%%%%%%%%%%%%%%%%%%%%
% 以上の設定には手を加えないこと - ここまで
%%%%%%%%%%%%%%%%%%%%%%%%%%%%%%%%%%%%%%%%%%%%%%%%%%%%%%%%%%%


\begin{document}
\maketitle
%
% 以下が問題と解答欄に関する記述
%
% 問題文には手を加えないこと．
% 解答は適切に編集して作成すること
%
%%%%%%%%%%%%%%%%%%%%%%%%%%%%%%%%%%%%%%%%%%%%%%%%%%%%%%%%%%%
%
% 問題1 
%
%%%%%%%%%%%%%%%%%%%%%%%%%%%%%%%%%%%%%%%%%%%%%%%%%%%%%%%%%%%
\section*{問題1}

演習室から利用可能な下記のディレクトリ
\begin{center}
\texttt{/cis/public/Literacy2022/}
\end{center}
がある．以下の問いに答えよ．
なお，各問いはこのディレクトリ配下にあるディレクトリやファイルを対象に考えよ．

%\begin{center}
 % \includegraphics[scale=0.5]{./prob1.png}
%\end{center}


%\begin{screen}
%\tt \$ bash fbuild.sh
%\end{screen}　

%%%%%%%%%%%%%%%%%%%%%%%%%%%%%%%%%%%%%%%%%%%%%%%%%%%%%%%%%%%
%
% 問題1-1
%
%%%%%%%%%%%%%%%%%%%%%%%%%%%%%%%%%%%%%%%%%%%%%%%%%%%%%%%%%%%
\subsection*{問題1-1}
\texttt{speechrec.txt}がどこにあるかを探したい．どのような操作をすればよいか．一行で示せ．ただし，検索結果は画面とホームディレクトリ直下の\texttt{result.txt}というファイルに出力することとする．


\answer{解答1-1}
%%%%%%%%%%%%%%%%%%%%%%%%%%%%%%%%%%%%%%%%%%%%%%%%%%%%%%%%%%%%
% 以下に解答を記述すること．
%%%%%%%%%%%%%%%%%%%%%%%%%%%%%%%%%%%%%%%%%%%%%%%%%%%%%%%%%%%
\begin{verbatim}
find /cis/public/Literacy2022/ -name "speechrec.txt" | tee ~/result.txt
\end{verbatim}

%%%%%%%%%%%%%%%%%%%%%%%%%%%%%%%%%%%%%%%%%%%%%%%%%%%%%%%%%%%%
%
% 問題1-2
%
%%%%%%%%%%%%%%%%%%%%%%%%%%%%%%%%%%%%%%%%%%%%%%%%%%%%%%%%%%%
\subsection*{問題1-2}
\texttt{speechrec.txt}をホームディレクトリにコピーするための操作を一行で示せ．なお，\texttt{speechrec.txt}は絶対パスで指定せよ．

\answer{解答1-2}
%%%%%%%%%%%%%%%%%%%%%%%%%%%%%%%%%%%%%%%%%%%%%%%%%%%%%%%%%%%
% 以下に解答を記述すること．
%%%%%%%%%%%%%%%%%%%%%%%%%%%%%%%%%%%%%%%%%%%%%%%%%%%%%%%%%%%
\begin{verbatim}
cp /cis/public/Literacy2022/2022_2_Literacy/dirA/speechrec.txt ~
\end{verbatim}


%%%%%%%%%%%%%%%%%%%%%%%%%%%%%%%%%%%%%%%%%%%%%%%%%%%%%%%%%%%
%
% 問題1-3
%
%%%%%%%%%%%%%%%%%%%%%%%%%%%%%%%%%%%%%%%%%%%%%%%%%%%%%%%%%%%
\subsection*{問題1-3}
\texttt{.wav}の文字列を含むファイルの名前とその文字列を含む行の数を表示したい．どのような操作をすればよいか．一行で示せ．なお，標準エラー出力は画面にもファイルにも出力しないようにせよ．

\answer{解答1-3}
%%%%%%%%%%%%%%%%%%%%%%%%%%%%%%%%%%%%%%%%%%%%%%%%%%%%%%%%%%%
% 以下に解答を記述すること．
%%%%%%%%%%%%%%%%%%%%%%%%%%%%%%%%%%%%%%%%%%%%%%%%%%%%%%%%%%%
\begin{verbatim}
find . -name "*" | xargs grep -c .wav 2> /dev/null
\end{verbatim}



%%%%%%%%%%%%%%%%%%%%%%%%%%%%%%%%%%%%%%%%%%%%%%%%%%%%%%%%%%%
%
% 問題1-4
%
%%%%%%%%%%%%%%%%%%%%%%%%%%%%%%%%%%%%%%%%%%%%%%%%%%%%%%%%%%%
\subsection*{問題1-4}
ホスト名\texttt{ie099}のコンピュータの\texttt{/tmp}ディレクトリにある\texttt{synthesis.txt}をホームディレクトリにコピーしたい．どのような操作をすればよいか．一行で示せ．

\answer{解答1-4}
%%%%%%%%%%%%%%%%%%%%%%%%%%%%%%%%%%%%%%%%%%%%%%%%%%%%%%%%%%%
% 以下に解答を記述すること．
%%%%%%%%%%%%%%%%%%%%%%%%%%%%%%%%%%%%%%%%%%%%%%%%%%%%%%%%%%%
\begin{verbatim}
scp ie099:/tmp/synthesis.txt ~
\end{verbatim}


%%%%%%%%%%%%%%%%%%%%%%%%%%%%%%%%%%%%%%%%%%%%%%%%%%%%%%%%%%%
%
% 問題1-5
%
%%%%%%%%%%%%%%%%%%%%%%%%%%%%%%%%%%%%%%%%%%%%%%%%%%%%%%%%%%%
\subsection*{問題1-5}
問題\texttt{1-2}でホームディレクトリにコピーした\texttt{speechrec.txt}の末尾に問題\texttt{1-3}のファイルの中身を追加したい．どのような操作をすればよいか．一行で示せ．
\answer{解答1-5}
%%%%%%%%%%%%%%%%%%%%%%%%%%%%%%%%%%%%%%%%%%%%%%%%%%%%%%%%%%%
% 以下に解答を記述すること．
%%%%%%%%%%%%%%%%%%%%%%%%%%%%%%%%%%%%%%%%%%%%%%%%%%%%%%%%%%%
\begin{verbatim}
find . -name "*" | xargs grep .wav | cat >> ~/speechrec.txt
\end{verbatim}

\newpage
%%%%%%%%%%%%%%%%%%%%%%%%%%%%%%%%%%%%%%%%%%%%%%%%%%%%%%%%%%%
%
% 問題2 
%
%%%%%%%%%%%%%%%%%%%%%%%%%%%%%%%%%%%%%%%%%%%%%%%%%%%%%%%%%%%
\input{user.inc}
\section*{問題2}
演習室のWebサーバ\url{www.cis.aichi-pu.ac.jp}で皆さんのウェブページを作成するものと仮定する．図2-1に示されているファイル\texttt{index.html}は，図2-2に示されているように閲覧されることを意図したHTMLドキュメントである．
\begin{center}
\IfFileExists{index-html.pdf}{\includegraphics[height=185mm]{index-html.pdf}}{\texttt{index-html.pdf}をカレントディレクトリに置いてください}
\\
図2-1  \texttt{index.html}ファイルの内容
\end{center}

ただし，出題の都合上，色は表示されず，また，【う】の箇所には別の内容が表示されるはずである．
\begin{center}
\IfFileExists{index.pdf}{\includegraphics[width=0.99\linewidth]{index.pdf}}{\texttt{index.pdf}をカレントディレクトリに置いてください}
\\
図2-2  \texttt{index.html}ファイルの内容
\end{center}
\newpage
以下の各問いに答えよ．URL, html, CSS, Javascriptなどの記述には半角文字を用いること．

\subsection*{問題2-1}
図2-1に示されている\texttt{index.html}が，
皆さん自身が作業をしている演習室のワークステーションの
カレントディレクトリにあるものとする．これを
\\
\hspace*{4zw}\URLa
\\
として公開するためにサーバ\texttt{www.cis.aichi-pu.ac.jp}にアップロードするための
コマンドを記せ．ただし，サーバには必要なディレクトリがすでに作成されているものとする．
\answer{解答2-1}
%%%%%%%%%%%%%%%%%%%%%%%%%%%%%%%%%%%%%%%%%%%%%%%%%%%%%%%%%%%
% 以下に解答を記述すること．
% varbatim環境を使うとよいだろう．
%%%%%%%%%%%%%%%%%%%%%%%%%%%%%%%%%%%%%%%%%%%%%%%%%%%%%%%%%%%
\begin{verbatim}
scp index.html www.cis.aichi-pu.ac.jp/~is221042/portfolio/report/
\end{verbatim}
\subsection*{問題2-2}
図2-1に示されている\texttt{index.html}に含まれるHTMLの文法の誤りを指摘せよ．
誤りの場所(○行目)と内容を記すこと．ただし，【あ】【い】【う】の部分は適切な内容に置き換えられているものとすること．

\answer{解答2-2}
%%%%%%%%%%%%%%%%%%%%%%%%%%%%%%%%%%%%%%%%%%%%%%%%%%%%%%%%%%%
% 以下に解答を記述すること．
%%%%%%%%%%%%%%%%%%%%%%%%%%%%%%%%%%%%%%%%%%%%%%%%%%%%%%%%%%%
(26行目)\verb%<h2>%に対し,\verb%</h3>%が対応している.
(42行目)\verb%<ol class="grade">%に対し(47行目)\verb%</ul>%が対応している.

\subsection*{問題2-3}
「成績評価」の中の4項目「1. 演習課題　30\%」，「2. 授業態度　30\%」，「3. 中間試験　20\%」，「4. 期末試験　20\%」がどの色で表示されるかを答えよ．
色の名前は\texttt{index.html}ファイルに記載されている通りに記すこと．

\answer{解答2-3}
%%%%%%%%%%%%%%%%%%%%%%%%%%%%%%%%%%%%%%%%%%%%%%%%%%%%%%%%%%%
% 以下に解答を記述すること．
%%%%%%%%%%%%%%%%%%%%%%%%%%%%%%%%%%%%%%%%%%%%%%%%%%%%%%%%%%%
1:teal,2:maroon,3:navy,4:purple


\subsection*{問題2-4}
「LaTeXによる文書作成」の文字を
\\
\hspace*{4zw}\URLb
\\
へのリンクとしたい．「LaTeXによる文書作成」の部分をどのように書き換えればよいか．相対パスで記すこと．
また，開始タグから終了タグまでのすべてを書くこと．
\answer{解答2-4}
%%%%%%%%%%%%%%%%%%%%%%%%%%%%%%%%%%%%%%%%%%%%%%%%%%%%%%%%%%%
% 以下に解答を記述すること．
% varbatim環境を使うとよいだろう．
%%%%%%%%%%%%%%%%%%%%%%%%%%%%%%%%%%%%%%%%%%%%%%%%%%%%%%%%%%%
\begin{verbatim}
<li>LaTeXによる文書作成 → <a href="https://www.cis.aichi-pu.ac
.jp/ ̃is221042/evidence/latex.html">LaTeXによる文書作成</a>
\end{verbatim}
\subsection*{問題2-5}
教員の名前の後ろにシルエット(\texttt{who.jpg})が表示されている．この画像をクリックするごとに
この場所に\texttt{iribe.jpg}, \texttt{ito.jpg}, \texttt{ohta.jpg}のいずれかがランダムに表示されるようにしたい．【あ】の部分に適切なスクリプトを記せ．
\answer{解答2-5}
%%%%%%%%%%%%%%%%%%%%%%%%%%%%%%%%%%%%%%%%%%%%%%%%%%%%%%%%%%%
% 以下に解答を記述すること．
% varbatim環境を使うとよいだろう．
%%%%%%%%%%%%%%%%%%%%%%%%%%%%%%%%%%%%%%%%%%%%%%%%%%%%%%%%%
\begin{verbatim}
  	let r;
        r = Math.floor(Math.random()*3);
        return 'photos[r]';
\end{verbatim}


\subsection*{問題2-6}
円グラフの画像について，画像が表示できないときに「成績基準」というテキストが表示されるようにするための記述を【い】に記せ．
\answer{解答2-6}
%%%%%%%%%%%%%%%%%%%%%%%%%%%%%%%%%%%%%%%%%%%%%%%%%%%%%%%%%%%
% 以下に解答を記述すること．
%%%%%%%%%%%%%%%%%%%%%%%%%%%%%%%%%%%%%%%%%%%%%%%%%%%%%%%%%%%



\subsection*{問題2-7}
50行目から54行目のスクリプトによって【う】の部分に表示される内容を記せ．
\answer{解答2-7}
%%%%%%%%%%%%%%%%%%%%%%%%%%%%%%%%%%%%%%%%%%%%%%%%%%%%%%%%%%%
% 以下に解答を記述すること．
%%%%%%%%%%%%%%%%%%%%%%%%%%%%%%%%%%%%%%%%%%%%%%%%%%%%%%%%%%%

``f(7416,3914)=''f(206,0)

\newpage
%%%%%%%%%%%%%%%%%%%%%%%%%%%%%%%%%%%%%%%%%%%%%%%%%%%%%%%%%%%
%
% 問題3
%
%%%%%%%%%%%%%%%%%%%%%%%%%%%%%%%%%%%%%%%%%%%%%%%%%%%%%%%%%%%
\section*{問題3}
以下の問いに答えよ．各問は演習室の\texttt{CentOS}環境で実施した場合を想定している．

\subsection*{問題3-1}
バックグラウンドで実行しているコマンド（実在しないがコマンド名は\texttt{commd}とする）が原因でコンピュータの動作が遅くなっている．そのため，そのコマンドのプロセスあるいはジョブを終了させたい，その手順とコマンド操作を説明せよ．
\answer{解答3-1}
%%%%%%%%%%%%%%%%%%%%%%%%%%%%%%%%%%%%%%%%%%%%%%%%%%%%%%%%%%%
% 以下に解答を記述すること．
%%%%%%%%%%%%%%%%%%%%%%%%%%%%%%%%%%%%%%%%%%%%%%%%%%%%%%%%%%%
kill -9 (ジョブ番号 \quad もしくは \quad プロセスID)で停止させる.
ジョブ番号は\$\quad jobsで調べる.プロセスIDは\$\quad ps\quad -e か\$\quad ps\quad -ef か\$\quad ps\quad -elyで調べると良い。
\subsection*{問題3-2}
ホームディレクトリの中身をリスト表示し，作成年の古い順に並び替えたい．そして，並び替えたリストをホームディレクトリ下の\texttt{list.txt}に出力したい．どのような操作をすればよいか．1行で示せ．
\answer{解答3-2}
%%%%%%%%%%%%%%%%%%%%%%%%%%%%%%%%%%%%%%%%%%%%%%%%%%%%%%%%%%%
% 以下に解答を記述すること．
%%%%%%%%%%%%%%%%%%%%%%%%%%%%%%%%%%%%%%%%%%%%%%%%%%%%%%%%%%%
\begin{verbatim}
ls -rlt ~/ > ~/list.txt
\end{verbatim}
\subsection*{問題3-3}
問題\texttt{3-2}で作成したファイルがホームディレクトリに存在するかどうかを確認するためのシェルスクリプト\texttt{check.sh}を作成せよ．
存在する場合は「存在する」と画面に表示し，存在しない場合はファイル\texttt{list.txt}をホームディレクトリに作成せよ．解答欄には\texttt{lstlisting}環境を用いてシェルスクリプトを記載すること．
\answer{解答3-3}
%%%%%%%%%%%%%%%%%%%%%%%%%%%%%%%%%%%%%%%%%%%%%%%%%%%%%%%%%%%
% 以下に解答を記述すること．
%%%%%%%%%%%%%%%%%%%%%%%%%%%%%%%%%%%%%%%%%%%%%%%%%%%%%%%%%%%
\begin{lstlisting}
#!/bin/sh

if [ -f ~/list.txt ] 
  then echo "存在する"
  else touch ~/list.txt
fi

\end{lstlisting}
\subsection*{問題3-4}
問題\texttt{3-3}で作成したシェルスクリプト \texttt{check.sh} を実行する方法を2通り示せ．
\answer{解答3-4}
%%%%%%%%%%%%%%%%%%%%%%%%%%%%%%%%%%%%%%%%%%%%%%%%%%%%%%%%%%%
% 以下に解答を記述すること．
%%%%%%%%%%%%%%%%%%%%%%%%%%%%%%%%%%%%%%%%%%%%%%%%%%%%%%%%%%%
sh ~/check.sh
~/check.sh (chmod 700 ~/test.sh 実行後)
\subsection*{問題3-5}
標準入力から数値を端末上に入力し，1からその数値までの和を出力するシェルスクリプト\texttt{sum.sh}を作成せよ．例えば，端末上に\texttt{100}を入力した場合，\texttt{1＋2＋3＋・・・＋98＋99＋100}を計算した結果（この場合\texttt{5050}）が画面に表示されるようにする．解答欄には \texttt{lstlisting}環境を用いてシェルスクリプトを記載すること
\answer{解答3-5}
%%%%%%%%%%%%%%%%%%%%%%%%%%%%%%%%%%%%%%%%%%%%%%%%%%%%%%%%%%%
% 以下に解答を記述すること．
%%%%%%%%%%%%%%%%%%%%%%%%%%%%%%%%%%%%%%%%%%%%%%%%%%%%%%%%%%%
\begin{lstlisting}
#!/bin/sh

read number

r=$(($number/2*($number+1)))
echo $r

\end{lstlisting}
\newpage
%%%%%%%%%%%%%%%%%%%%%%%%%%%%%%%%%%%%%%%%%%%%%%%%%%%%%%%%%%%
%
% 問題4
%
%%%%%%%%%%%%%%%%%%%%%%%%%%%%%%%%%%%%%%%%%%%%%%%%%%%%%%%%%%%
\section*{問題4}
問題~4-1, 問題~4-2それぞれに示されるBeamerスライドの画像に
対応する{\LaTeX}~ソースコードを推測・再構成し，解答せよ．
また，その{\LaTeX}ソースコードから出力できるPDFを示せ．
ただし，以下の点に注意せよ：
\begin{itemize}
\item
  Beamerクラスのオプションで，アスペクト比は 16:9，フォントサイズは
  12ポイントを指定すること；
\item 
  スライドの各種テーマや改行のタイミングなどは不問だが，内容の
  ポイントは逃さないようにすること
  （ヒント：それぞれ 8up，4up での出力だが，重ね描き機能がON/OFFの違い
  がある．ハンドアウトオプションをつけるかどうかで…）；
\item
  墨付き括弧（例：【今日の年月日と曜日】）の部分は，
  各自にとって適切な情報に置き換えること；
% \item
%   \verb|\textcolor{【色の名前】}{【色をつけたい文字列】}|~コマンド
%   を利用すれば，文字列に色をつけることができる；
\item
  Beamerスライドでは，プリアンブルに
  \verb|\setbeamertemplate{caption}[numbered]| を書いておけば，
  図や表のカウンター（自動でつく番号）を明示できる；
% \item
%   Beamerスライドでは，\verb|thebibliography|~環境の内側に
%   \verb|\beamertemplatetextbibitems|を書いておけば，参考文献リスト
%   のカウンター（自動でつく番号）を明示できる；
\item
  Beamerスライドにおいて \verb|verb|~コマンドを使いたい場合，
  \verb|frame|~環境に \verb|fragile|~オプション を書いておく
  （つまり，\verb|\begin{frame}|を\verb|\begin{frame}[fragile]|
  に置き換える）こと．さもなくば，コンパイルがエラー終了してしまう．
\end{itemize}
\pagebreak

\subsection*{問題4-1}
\begin{center}
  \includegraphics[width=\textwidth]{./prob4-1.pdf}
\end{center}
\pagebreak

\answer{解答4-1　{\LaTeX}~ソースコード}
%%%%%%%%%%%%%%%%%%%%%%%%%%%%%%%%%%%%%%%%%%%%%%%%%%%%%%%%%%%
% LaTeXコンパイルで，ソースリストが表示されるようにしてください
%%%%%%%%%%%%%%%%%%%%%%%%%%%%%%%%%%%%%%%%%%%%%%%%%%%%%%%%%%%
%
% uncomment the following line(s) if necessary
% 
% \begin{listing}[!h]
%   \caption{ans4-1.tex}
%   \label{code:ans4-1}
%   \inputminted[%
%   linenos,numbersep=5pt,breaklines,%
%   frame=lines,
%   fontsize=\footnotesize]{latex}{./ans4-1.tex}
% \end{listing}
%
% or
% 
%\lstinputlisting[caption=ans4-1.tex,language=TeX,label=ans4-1]{./my-ans4-1.tex}
%
% or
% 
% \begin{verbatim}
% ...
% \end{verbatim}
%
\pagebreak
%
\answer{解答4-1　PDF出力}
%%%%%%%%%%%%%%%%%%%%%%%%%%%%%%%%%%%%%%%%%%%%%%%%%%%%%%%%%%%
% 出力結果のPDFを取り込む
%%%%%%%%%%%%%%%%%%%%%%%%%%%%%%%%%%%%%%%%%%%%%%%%%%%%%%%%%%%
% \begin{center}
%   \includegraphics[width=.95\textwidth]{~/my-ans4-1.pdf}
% \end{center}
\pagebreak

\subsection*{問題4-2}
\begin{center}
  \includegraphics[width=\textwidth]{./prob4-2.pdf}
\end{center}
\pagebreak

\answer{解答4-2　{\LaTeX}~ソースコード}
%%%%%%%%%%%%%%%%%%%%%%%%%%%%%%%%%%%%%%%%%%%%%%%%%%%%%%%%%%%
% LaTeXコンパイルで，ソースリストが表示されるようにしてください
%%%%%%%%%%%%%%%%%%%%%%%%%%%%%%%%%%%%%%%%%%%%%%%%%%%%%%%%%%%
%
% uncomment the following line(s) if necessary
% 
% \begin{listing}[!h]
%   \caption{ans4-2.tex}
%   \label{code:ans4-2}
%   \inputminted[%
%   linenos,numbersep=5pt,breaklines,%
%   frame=lines,
%   fontsize=\footnotesize]{latex}{./ans4-2.tex}
% \end{listing}
%
% or
% 
%\lstinputlisting[caption=ans4-2.tex,language=TeX,label=ans4-2]{./my-ans4-2.tex}
%
% or
% 
% \begin{verbatim}
% ...
% \end{verbatim}
%
\pagebreak
%
\answer{解答4-2　PDF出力}
%%%%%%%%%%%%%%%%%%%%%%%%%%%%%%%%%%%%%%%%%%%%%%%%%%%%%%%%%%%
% 出力結果のPDFを取り込む
%%%%%%%%%%%%%%%%%%%%%%%%%%%%%%%%%%%%%%%%%%%%%%%%%%%%%%%%%%%
% \begin{center}
%   \includegraphics[width=.95\textwidth]{./my-ans4-2.pdf}
% \end{center}


\newpage
%%%%%%%%%%%%%%%%%%%%%%%%%%%%%%%%%%%%%%%%%%%%%%%%%%%%%%%%%%%
%
% 問題5-1
%
%%%%%%%%%%%%%%%%%%%%%%%%%%%%%%%%%%%%%%%%%%%%%%%%%%%%%%%%%%%
\section*{問題5}
\subsection*{問題5-1}
以下の計算問題を解け．3.14159や0.333333のような数値解ではなく，
$\pi$や$\frac{1}{3}$のような厳密解を求めること．
Mathematicaを用いた場合は，用いたコマンドも記すこと．
問題5-1-1, 5-1-2については，問題とともに結果を記せ．途中の計算は不要である．
\begin{quote}
例　問題：$m$ を皆さんの誕生月とするときに，$m+2$ を計算せよ．\\
　　解答(5月生まれの場合)\\
　　　　Mathematicaコマンド \verb@5+2@\\
　　　　解 $5+2=7$
\end{quote}
\subsubsection*{問題5-1-1}
次の和を計算せよ．
\begin{center}
\IfFileExists{./sum.pdf}{
\includegraphics[height=8.5ex]{sum.pdf}}{\texttt{sum.pdf}をカレントディレクトリに置いてください}
\end{center}
\answer{解答5-1-1}
%%%%%%%%%%%%%%%%%%%%%%%%%%%%%%%%%%%%%%%%%%%%%%%%%%%%%%%%%%%
% 以下に解答を記述すること．
%%%%%%%%%%%%%%%%%%%%%%%%%%%%%%%%%%%%%%%%%%%%%%%%%%%%%%%%%%%
\begin{verbatim}
コマンド
Sum[((3/11)^n)*n*(n - 1)*(n - 2), {n, \[Infinity]}]
解
891/2048

\[] は　エスケープ　→　inf →　エスケープ
\end{verbatim}
\subsubsection*{問題5-1-2}
次の積分を計算せよ．Mathematicaは分母に根号が残った解を出してくるかもしれないが，有理化する必要はない．
\begin{center}
\IfFileExists{./integ.pdf}{\includegraphics[height=7.5ex]{integ.pdf}}{\texttt{integ.pdf}をカレントディレクトリに置いてください}
\end{center}
\answer{解答5-1-2}
%%%%%%%%%%%%%%%%%%%%%%%%%%%%%%%%%%%%%%%%%%%%%%%%%%%%%%%%%%%
% 以下に解答を記述すること．
%%%%%%%%%%%%%%%%%%%%%%%%%%%%%%%%%%%%%%%%%%%%%%%%%%%%%%%%%%%
\begin{verbatim}
コマンド
Integrate[1/((4 x^2 + 22 x + 31)^6), {x, -\[Infinity], \[Infinity]}]
解
(28 \[Pi])/(27 Sqrt[3])
\end{verbatim}
\subsubsection*{問題5-1-3}
連立1次方程式 $ax=b$ を解け．ただし，$a$, $b$ は以下で与えられるものとする．解の個数にも気を付けること．
\begin{center}
\IfFileExists{./matrix.pdf}{\includegraphics[height=15ex]{matrix.pdf}}{\texttt{matrix.pdf}をカレントディレクトリに置いてください}
\end{center}
\answer{解答5-1-3}
%%%%%%%%%%%%%%%%%%%%%%%%%%%%%%%%%%%%%%%%%%%%%%%%%%%%%%%%%%%
% 以下に解答を記述すること．
%%%%%%%%%%%%%%%%%%%%%%%%%%%%%%%%%%%%%%%%%%%%%%%%%%%%%%%%%%%
\begin{verbatim}
コマンド
a = {{25, -18, -46, -6, -14}, {4, -2, -8, 0, -1}, {-2, 2, 3, 1, 
   2}, {12, -9, -22, -3, -7}, {-48, 33, 89, 10, 25}}
b = {{-24}, {0}, {5}, {-13}, {40}}
LinearSolve[a, b]
解
{{-2}, {4}, {-2}, {-1}, {0}}
\end{verbatim}
\newpage
%%%%%%%%%%%%%%%%%%%%%%%%%%%%%%%%%%%%%%%%%%%%%%%%%%%%%%%%%%%
%
% 問題5-2
%
%%%%%%%%%%%%%%%%%%%%%%%%%%%%%%%%%%%%%%%%%%%%%%%%%%%%%%%%%%%
\subsection*{問題5-2}
MATLABを用いて以下の各問いに答えよ．作成したコードも示すこと．

なお，別途配布する\texttt{plot-sample.pdf}にグラフの出力イメージが示されている．
\begin{center}
\framebox{\begin{minipage}{0.95\linewidth}【補足】グラフの出力について
\begin{itemize}\item コマンドを用いてEPSファイルに出力する場合には，\texttt{-deps}の代わりに\texttt{-depsc}を用いてカラーで出力すること．
\item PDFファイルに出力した場合には，コマンド，GUIにかかわらず，図の周りに大きな余白があるので，（MATLABではなく端末から)\texttt{pdfcrop}コマンドを用いて削除すること．例えば，PDFファイルの名前を\texttt{plot1.pdf}とするとき，
\\ \hspace*{4zw}
\texttt{pdfcrop plot1.pdf}
\\
によって余白が消去されたファイル\texttt{plot1-crop.pdf}が作られる。\end{itemize}\end{minipage}}
\end{center}

\subsubsection*{問題5-2-1}
複素数 $\omega=\OMEGA$ に対して，
複素数平面上に$\omega^n$ ($n=0,1,\cdots,32$)を赤色のサイズ10のマーカー`+'で示すとともに，曲線$\omega^t$ ($0\leq t\leq 32$)を青色の実線で示せ．また，$x$軸と$y$軸も表示すること．



\answer{解答5-2-1　結果（グラフ）}
%%%%%%%%%%%%%%%%%%%%%%%%%%%%%%%%%%%%%%%%%%%%%%%%%%%%%%%%%%%
% 結果のグラフの画像を取り込む
%%%%%%%%%%%%%%%%%%%%%%%%%%%%%%%%%%%%%%%%%%%%%%%%%%%%%%%%%%%
% \begin{center}
%   \includegraphics[width=.8\textwidth]{./prob5-2-1.pdf}
% \end{center}

\answer{解答5-2-1　MATLABコード}
%%%%%%%%%%%%%%%%%%%%%%%%%%%%%%%%%%%%%%%%%%%%%%%%%%%%%%%%%%%
% LaTeXコンパイルで，MATLABコードが表示されるようにしてください
%%%%%%%%%%%%%%%%%%%%%%%%%%%%%%%%%%%%%%%%%%%%%%%%%%%%%%%%%%%
%
% 必要に応じて以下の行のコメントを削除して利用できます
% 
%\begin{lstlisting}[caption=prob5-2-1,language=Matlab,label=prob5-2-1]
% ここにコードを書く
%
%\end{lstlisting}
%
% または
%
%\begin{listing}[!ht]
%  \caption{prob5-2-1}
%  \label{prob5-2-1}
% \begin{minted}[linenos,numbersep=5pt,breaklines, frame=lines, fontsize=\footnotesize]{matlab}
% ここにコードを書く
%
% \end{minted}
%\end{listing}

\subsubsection*{問題5-2-2}
配布されている\verb@temp.dat@は\CITY{}市の\YEAR{}年の毎日の気温のデータ\cite{jma}を並べたもので，これを\verb@load@コマンドで変数\verb@a@に読み込むと30行12列の行列が得られる．\verb@a@の各列は(おおよそ)各月の気温のデータである(正確には平年では第1列は1月1日から1月30日，第2列は1月31日から3月1日，第3列は3月2日から3月31日，うるう年では第1列は1月1日から1月30日，第2列は1月31日から2月29日，第3列は3月1日から3月30日などであるが，気にしないことにする)．

この行列\texttt{a}に対して，\texttt{boxplot(a)}を実行すると，各月の気温の分布が箱ひげ図で描かれる．この図に，\texttt{a}の各列の平均値として得られる，各月の気温の平均値を表す折れ線グラフを付け加えよ．各月の気温の平均値を\texttt{t1}, \texttt{t2}, ...とするとき，点\texttt{(1,t1)}, \texttt{(2, t2)}, ...を結べばよい．
\begin{itemize}
\item 横軸には「Month」縦軸には「Temperature [$^\circ$C]」というラベルをつけよ．記号$^\circ$CはMATLABでも\LaTeX{}流の \verb@^\circ{}C@ で入力できる(前後の \verb@$@ は不要)。
\item 折れ線グラフは，緑色で正方形のマーカーとそれらを結ぶ実線で描け．マーカーのサイズは10に設定せよ．
\item
  \verb@set(gca, 'FontSize', 12)@コマンドを用いて
  フォントサイズを 12~ポイントにせよ．
\item
  グラフタイトルは何もつけないこと．
\item
  \texttt{load}コマンドでデータを読み込むところからグラフが完成するまでのコードを示すこと．
\end{itemize}

\answer{解答5-2-2　結果（グラフ）}
%%%%%%%%%%%%%%%%%%%%%%%%%%%%%%%%%%%%%%%%%%%%%%%%%%%%%%%%%%%
% 結果のグラフの画像を取り込む
%%%%%%%%%%%%%%%%%%%%%%%%%%%%%%%%%%%%%%%%%%%%%%%%%%%%%%%%%%%
% \begin{center}
%   \includegraphics[width=.8\textwidth]{./prob5-2-2.pdf}
% \end{center}

\answer{解答5-2-2　MATLABコード}
%%%%%%%%%%%%%%%%%%%%%%%%%%%%%%%%%%%%%%%%%%%%%%%%%%%%%%%%%%%
% LaTeXコンパイルで，MATLABコードが表示されるようにしてください
%%%%%%%%%%%%%%%%%%%%%%%%%%%%%%%%%%%%%%%%%%%%%%%%%%%%%%%%%%%
%
% 必要に応じて以下の行のコメントを削除して利用できます
% 
%\begin{lstlisting}[caption=prob5-2-2,language=Matlab,label=prob5-2-2]
% ここにコードを書く
%
%\end{lstlisting}
%
% または
%
%\begin{listing}[!ht]
%  \caption{prob5-2-2}
%  \label{prob5-2-2}
% \begin{minted}[linenos,numbersep=5pt,breaklines, frame=lines, fontsize=\footnotesize]{matlab}
% ここにコードを書く
%
% \end{minted}
%\end{listing}


%%%%%%%%%%%%%%%
\begin{thebibliography}{9}
\bibitem{jma}気象庁，
過去の気象データ検索，
\url{https://www.data.jma.go.jp/obd/stats/etrn/index.php}，2022年7月28日閲覧
\end{thebibliography}
\end{document}
